\section{训练监控与排障要点}
\subsection{从奖励曲线反推系统问题}
在真正运行 Agentic RL 任务时，研究者往往不会先看到“算法错了”，而是先看到 reward curve 长时间不抬升、方差异常大或者 rollout 质量忽高忽低。把这些现象拆解开，通常能更快定位到问题究竟来自 reward、环境还是采样配置。

\begin{table}[H]
\centering
\begin{tabular}{p{3.5cm}p{5cm}p{4.5cm}}
\toprule
观测现象 & 常见原因 & 首先检查什么 \\
\midrule
reward 长期不涨 & reward 定义太稀疏或 credit assignment 太弱 & 是否需要 shaping 或更细粒度反馈 \\
loss 波动很大 & batch 太小或 advantage 方差过高 & rollout 长度、归一化与 clip 配置 \\
策略突然崩掉 & environment 不稳定或更新过猛 & 采样日志、异常 episode、KL 漂移 \\
\bottomrule
\end{tabular}
\caption{Agentic RL 训练里最常见的三类异常信号}
\end{table}

\begin{knowledgebox}{为什么短视频也要强调监控}
即便这一讲只介绍 PG loss 的核心组件，真正把它用到 veRL 时，工程反馈仍然是第一位的。原因很简单：策略梯度类方法对采样分布十分敏感，很多问题只有在 reward、KL、episode length、成功率同时观察时才看得出来。
\end{knowledgebox}

\subsection{本章小结}
PG loss 的理论理解必须和训练监控结合起来：只有把 reward、rollout 质量和策略漂移一起看，才能判断当前更新到底是在学习，还是只是在噪声里震荡。

\newpage
